\subsection{Indentation Style}\label{sec:IndentationStyle}
The indentation\index{indentation} of code blocks, based on their context, increases the readability of the code.

Obviously, \bdq{readability} is a relative term, depending from one's habits. So some people are more familiar with the Kernighan-Ritchie (K\&R) indentation style\index{Kernighan-Ritchie style}\index{K\&R style|see{Kernighan-Ritchie style}}

\begin{lstlisting}
public final void myMethod (){
    if (flag) {
        …
    }
}   //  myMethod()
\end{lstlisting}

as it is also shown in the \ccjpl, other prefer to have the curly braces on a line of their own (sometimes referred to as GNU or BSD style\index{GNU style}\index{BSD style}\footnote{All, BSD, GNU and K\&R styles, will define more than just how to place the curly braces. In addition, all three are originally part of code conventions for programming in C and/or \mbox{C++} that cannot applied to Java without modifications. See chapter \tqfullvref{sec:OtherProgrammingLanguages} on this topic.}):

\begin{lstlisting}
public final void myMethod()
{
    if( flag )
    {
    	…
    }
}   //  myMethod()
\end{lstlisting}

Both styles will work fine, but they look different. So as a lot of other guidelines from a code conventions document, too, this is a matter of taste. But as conformity supports readability, a main purpose of these guidelines is to ensure that all source code looks similar, and we have to determine the style to use.

I prefer the so-called GNU or BSD style\index{indentation!style} and recommend it to you as well, because it makes it easier to detect the begin and the end of a code block than the Kernighan-Ritchie style. It demands to have the opening and closing curly braces on a line of their own, like below. This let the corresponding opening and closing braces appear on the same column:

\begin{lstlisting}
public final class MyClass
{
    public final void myMethod()
    {
        if( flag )
        {
            /* do something */
        }
        
        for( final var i = 0; i < max; ++i )
        {
            /* do something else */
        } 
    }   //  myMethod()
}   //  class MyClass
\end{lstlisting}